\documentclass[10pt,a4paper]{amsart}
\usepackage[margin=19mm]{geometry}
\usepackage{amsmath,amssymb,mathtools}
\usepackage[scr=rsfs,cal=euler]{mathalfa}
\usepackage{xfrac}
\usepackage[unicode,hidelinks,bookmarks=false]{hyperref}
\newcommand{\ie}{i.e.\ }
\newcommand{\cf}{cf.\ }
\newcommand{\resp}{respectively\ }
\newcommand{\Subsection}{\subsection}
\providecommand{\nobreakdash}{\nobreak\hbox{-}}
\setlength{\emergencystretch}{2em}
\multlinegap0pt
\usepackage{float}
\usepackage{amssymb}
\usepackage{tikz}
\usepackage{dsfont}
\usepackage{bm}
\usepackage{enumerate}
\newtheorem{proposition}{Proposition}
\newtheorem{remark}{Remark}
\newtheorem{theorem}{Theorem}
\title{Time-periodic solutions for viscous fluids interacting with nonlinear Koiter plates}
\author{Claudiu M\^indril\u{a}}
\date{Source: arXiv:2605.19051v1 (CC BY 4.0)}
\begin{document}
\maketitle

\noindent\emph{Source and licence.} Time-periodic solutions for viscous fluids interacting with nonlinear Koiter plates. Version arXiv:2605.19051v1 (CC BY 4.0), distributed by arXiv under the Creative Commons Attribution 4.0 International licence, http://creativecommons.org/licenses/by/4.0/, as recorded in the arXiv licence metadata for this exact version and shown on the abstract page https://arxiv.org/abs/2605.19051v1. Full licence notice: \texttt{licenses/arxiv-2605.19051-CC-BY-4.0.txt}.\par\smallskip

\noindent\textit{Editorial scope.} Counted unit: the article's proposition proposition (source0.tex lines 1229--1231). The article's complete original proof of it is delivered with it (source0.tex lines 1232--1275, one \texttt{proof} environment of 44 lines). No mathematics is written, repaired or re-derived: the statement and the proof are the article's own lines, in the article's own notation, and no retained line is substituted or re-flowed.  Retained uncounted as the source-local context the proof needs: lines 1086--1092; lines 1122--1124; lines 1139--1153; lines 1684--1687.  Nothing was omitted from any retained span, so the package carries no omission record.  No retained line contains a citation, so the wrapper carries no bibliography and no bibliography entry.  No retained line contains a figure environment, an image-inclusion command or drawing code, so no image is delivered and the package declares no asset dependency.  Editorial formatting only, in addition: an AMS \texttt{amsart} preamble that restates the article's own declarations and macros used by the retained lines, namely the environments \texttt{proposition}, \texttt{remark}, \texttt{theorem} on separate counters, together with the standard packages those lines need.  The retained environments print in this document as Proposition 1, Remark 1, Theorem 1 in order of appearance, whereas the article prints the same items with the numbering of its own full document.  The complete native source of the article as distributed in the arXiv e-print for this exact version is bundled unchanged as \texttt{sources/200889/source0.tex}.
\begin{equation}\label{eqn:xk}
\begin{aligned}\int_{\Omega_{\delta\left(t\right)}}\partial_{t}\mathbf{u}_{n}\cdot\mathbf{X}_{k}dx+\frac{1}{2}\int_{\omega}\partial_{t}\eta_{n}\partial_{t}\delta X_{k}dA+\int_{\Omega_{\delta\left(t\right)}}\nabla\mathbf{u}_{n}\cdot\nabla\mathbf{X}_{k}dx & +\\
\frac{1}{2}\int_{\Omega_{\delta\left(t\right)}}\left(\mathbf{u}_{n}\cdot\nabla\right)\mathbf{u}_{n}\cdot\mathbf{X}_{k}dx-\frac{1}{2}\int_{\Omega_{\delta\left(t\right)}}\left(\mathbf{u}_{n}\cdot\nabla\right)\mathbf{X}_{k}\cdot\mathbf{u}_{n}dx & +\\
\int_{\omega}\partial_{tt}\eta_{n}X_{k}dA+\left\langle K^{\prime}\left(\eta_{n}\right),X_{k}\right\rangle  & =\\
\int_{\Omega_{\delta\left(t\right)}}\mathbf{f}\cdot\mathbf{X}_{k}dx+\int_{\omega}gX_{k}dA\quad\forall t\in I
\end{aligned}
\end{equation}

\begin{equation}\label{eqn:cons-gronwall}
\sup_{t\in\left[0,T\right]}E_{n}\left(t\right)<\infty 
\end{equation}

\begin{proposition}\label{prop: invariant-ball} 
Suppose that \[E_{n} (T) \ge E_{n} (0).\]
Then
there exists a constant    
\begin{equation}\label{eqn:R-chap02}
R=c\left(\texttt{data}\right)\cdot\left(C\left(\mathbf{f},g\right)^{2}+C\left(\mathbf{f},g\right)+m^{2}\right)>0
\end{equation}
for which
 \[E_{n} \left(0 \right) \le E_{n} \left(T \right) \le R.\]
By integrating in time from $0$ to $t\in I$ we have
\begin{equation}\label{eqn:chap02-bound-E-n}
\sup_{n\ge1}\sup_{t\in I}E_{n}\left(t\right)\lesssim R+C(\mathbf{f},g)<\infty.
\end{equation}

\end{proposition}

\begin{theorem}\label{thm:Schaeffer}
Let $X$ be a Banach space and   $A :  X \mapsto X$ be a continuous and compact mapping such that the set \[\left\{ x\in X:x=\lambda Ax\ \text{for some}\ \lambda\in\left[0,1\right]\right\} \]
is bounded. Then $A$ has at least one fixed point.
\end{theorem}

\begin{proposition}
The mapping $ \mathcal{T}:X \mapsto X$ admits at least one fixed point.
\end{proposition}

\begin{proof}
We shall prove that $\mathcal{T}$ satisfies the conditions of the  Leray-Schauder Theorem~\ref{thm:Schaeffer}. Note that $\mathcal{T}$ is well defined, recalling the discussion around \eqref{eqn:cons-gronwall}.

The compactness of $\mathcal{T}$ is due to the Arzela-Ascoli theorem and the fact that
\[
\mathcal{T}\left(X\right)\subset C^{2}\left(I\right)\hookrightarrow\hookrightarrow C^{1}\left(I\right).
\] 

That $\mathcal{T}$ is continuous follows from the well-posedness of \eqref{eqn:xk}. 
Indeed, let us consider a sequence $x_{k}\to x$ in $X$ and prove that $\mathcal{T}x_{k}\to\mathcal{T}x$.
The key observation is that if $\mathcal{T}x_{k_{n}}$ is an arbitrary convergent subsequence of $\mathcal{T}x_{k}$, then its limit equals $\mathcal{T}x$. This is because the limit would correspond to a solution of \eqref{eqn:xk} with the initial values and the domain corresponding to $x$. From the well-posedness of the system, this is exactly $\mathcal{T}x$.

To finish the argument, suppose there exists $\varepsilon_0>0$ and a subsequence of $\mathcal{T}x_{k}$ (not relabeled) for which 
\begin{equation}\label{eqn:contrad}
\left\Vert \mathcal{T}x_{k}-\mathcal{T}x\right\Vert _{X}\ge\varepsilon_{0}>0.
\end{equation}
But from \eqref{eqn:cons-gronwall} and \eqref{eqn:xk} it follows that $\mathcal{T}x_{k}$ is bounded in $C^{2}\left(I\right)$. Hence, it has a strongly convergent subsequence. Its limit, however, has to equal $\mathcal{T}x$, as proved before. This contradicts \eqref{eqn:contrad} and proves that $\mathcal{T}$ is continuous.

The most delicate part consists in proving that the set 
\[LS:=\left\{ x\in X:x=\lambda\mathcal{T}x\ \text{for some}\ \lambda\in\left[0,1\right]\right\} 
\] is uniformly bounded. 
To this end, note that in accordance with Proposition~\ref{prop: invariant-ball} it suffices to prove that for any $\mathbf{a}_{n}\in LS$ it holds that $E_{n}\left(T\right)\ge E_{n}\left(0\right)$. 
Since $\mathbf{a}_{n} \in LS$ it follows that 
\begin{equation}\label{eqn:cond-ls}
\left(\mathbf{a}_{n},\mathbf{b}_{n}\left(0\right),\mathbf{b}_{n}^{\prime}\left(0\right)\right)=\lambda\left(\mathbf{b}_{n},\mathbf{b}_{n}\left(T\right),\mathbf{b}_{n}^{\prime}\left(T\right)\right),\quad\text{for}\ \lambda\in\left[0,1\right].
\end{equation}
The case $\lambda=0$ is trivial, so we assume $0<\lambda \le 1$.
Then by  \eqref{eqn:cond-ls} we can compute
\begin{equation}\begin{aligned}E_{n}\left(T\right)-E_{n}\left(0\right)= & \left(\int_{\Omega_{\delta\left(T\right)}}\left|\mathbf{u}_{n}\left(T\right)\right|^{2}dx-\int_{\Omega_{\delta\left(0\right)}}\left|\mathbf{u}_{n}\left(0\right)\right|^{2}dx\right)+\\
 & \int_{\omega}\left(\partial_{t}\eta\left(T\right)\right)^{2}-\left(\partial_{t}\eta\left(0\right)\right)^{2}dA+\\
 & \left(K\left(\frac{1}{\lambda}\eta\left(0\right)\right)-K\left(\eta\left(0\right)\right)\right)\\
\ge & \left(\frac{1}{\lambda^{2}}-1\right)\left[\int_{\Omega_{\delta\left(0\right)}}\left|\mathbf{u}_{n}\left(0\right)\right|^{2}dx+\int_{\omega}\left(\partial_{t}\eta\left(0\right)\right)^{2}dA\right]+\\
 & \left(\frac{1}{\lambda^{2}}-1\right)K\left(\eta\left(0\right)\right)\\
\ge & 0.
\end{aligned}
\end{equation}

\begin{remark}
Note that comparing integrals on $\Omega_{\delta(0)}$ and $\Omega_{\delta(T)}$ is non-trivial when the domains differ. This is precisely why we require $\delta(0)=\delta(T)$, which is ensured by the operator $\mathcal{P}_{\varepsilon}$ above.
\end{remark}

This now places us in the context of Proposition~\ref{prop: invariant-ball}, so the set $LS$ is indeed uniformly bounded. With this, the conditions of the Leray-Schauder fixed-point theorem have been fulfilled and thus there exists a  fixed point of $\mathcal{T}$, that is an $\mathbf{a}_{n} \in X$ with 
\[\mathbf{a}_{n}=\mathbf{b}_{n},\quad\mathbf{b}_{n}\left(0\right)=\mathbf{b}_{n}\left(T\right),\quad\mathbf{b}_{n}^{\prime}\left(0\right)=\mathbf{b}_{n}^{\prime}\left(T\right).\]
\end{proof}

\end{document}
